\section{Introduction}

Interferometric observations of gravity darkening in stars
\citep[e.g.,][]{2018MNRAS.480.1263H,2018A&A...619A.167D} and the
recent progress of intensity interferometry
\citep{2025ApJ...995..191A} have created interest in basis functions
for stellar interferometry.

\cite{2026PASP..138e4504D} developed such a basis set for a rotating
star.  In this work we will re-derive their basis set with a
refinement that makes it closed-form, and extend it both
perturbatively and numerically to the case of a triaxial star, which
is flattened by spin and stretched by tide in along different axes.
This is done using symbolic computing and AI assistance.

A recent essay by \citet{tao2026} identifies five stages of a new
result in mathematics.
\begin{itemize}
\item \textit{Generation} of new theorem proofs.  Here AI has brought
  revolutionary changes.
\item \textit{Verification} of the proofs.  Symbolic computing had
  made it possible to keep up with the speed of generation.
\item \textit{Exposition} for human readers.  AI has not (yet) helped
  here, and in fact has made things worse --- the problem of AI slop.
\item \textit{Publication} in the research literature.
\item \textit{Canonicalization}, or absorption into textbooks.
\end{itemize}
Here would like to study how the first four stages translate into
astrophysics.

